%!TEX root = ../template.tex
%%%%%%%%%%%%%%%%%%%%%%%%%%%%%%%%%%%%%%%%%%%%%%%%%%%%%%%%%%%%%%%%%%%%
%% sdg-statement-pt.tex
%% NOVA thesis document file
%%
%% Statement printed below the SDG icons (optional, Portuguese)
%%%%%%%%%%%%%%%%%%%%%%%%%%%%%%%%%%%%%%%%%%%%%%%%%%%%%%%%%%%%%%%%%%%%

\typeout{NT FILE sdg-statement-pt.tex}%

% Enable with:
%     \ntaddfile{sdgstatement}[pt]{sdg-statement-pt}
% in 0-Config/4_files.tex, right below the abstracts.
%
% This particular text explains why the novathesis template itself
% contributes to SDGs 4, 9, 10 and 17 (set in 0-Config/1_novathesis.tex).
% If you change print/sdgs/list for your own document, rewrite this text
% to justify your own selection instead.

Enquanto modelo gratuito e de código aberto, o próprio \novathesis{} contribui para vários destes objetivos: reduz uma barreira real entre os estudantes e a conclusão do seu grau académico, em dezenas de escolas e vinte línguas (ODS~4, Educação de Qualidade); sendo gratuito e aberto, elimina o custo de ferramentas ou serviços de formatação pagos, alargando esse acesso a estudantes que de outra forma não o poderiam pagar (ODS~10, Redução das Desigualdades); é uma infraestrutura aberta e reutilizável sobre a qual outras instituições constroem as suas próprias configurações (ODS~9, Indústria, Inovação e Infraestruturas); e é mantido através de contribuições públicas, um programa de mantenedores por escola, e discussão aberta da comunidade (ODS~17, Parcerias para a Implementação dos Objetivos).
